\documentclass[11pt]{article}
\usepackage[T1]{fontenc}
\usepackage[utf8]{inputenc}
\usepackage{lmodern}
\usepackage[margin=1in]{geometry}
\usepackage{amsmath,amssymb,amsthm,mathtools}
\usepackage{microtype}
\usepackage{stmaryrd}
\usepackage{booktabs,longtable,array}
\usepackage{xcolor}
\usepackage{listings}
\usepackage{enumitem}
\usepackage{xurl}
\usepackage{hyperref}
\hypersetup{colorlinks=true,linkcolor=blue!45!black,citecolor=blue!45!black,urlcolor=blue!45!black,
 pdftitle={A Constructive Lean Contender Beyond Contender 5},pdfauthor={OpenAI; prepared for Vladimir Reshetnikov}}
\usepackage{fancyhdr}
\pagestyle{fancy}\fancyhf{}
\fancyhead[L]{\small Constructive large numbers in Lean}
\fancyhead[R]{\small Research proposal, 12 October 2026}
\fancyfoot[C]{\thepage}
\setlength{\headheight}{14pt}
\setlength{\parskip}{4pt}
\setlength{\parindent}{0pt}
\setlist{itemsep=2pt,topsep=4pt}
\emergencystretch=2em
\newtheorem{theorem}{Theorem}[section]
\newtheorem{lemma}[theorem]{Lemma}
\newtheorem{proposition}[theorem]{Proposition}
\newtheorem{corollary}[theorem]{Corollary}
\theoremstyle{definition}
\newtheorem{definition}[theorem]{Definition}
\newtheorem{remark}[theorem]{Remark}
\newcommand{\N}{\mathbb N}
\newcommand{\W}{\mathcal W}
\newcommand{\Sup}{\mathsf{sup}}
\newcommand{\leaf}{\mathbf 0}
\newcommand{\tyN}{\mathsf N}
\newcommand{\tyW}{\mathsf W}
\newcommand{\val}{\operatorname{val}}
\newcommand{\depth}{\operatorname{depth}}
\newcommand{\td}{\operatorname{td}}
\newcommand{\bits}{\operatorname{bits}}
\newcommand{\Qte}{\mathsf{quote}}
\newcommand{\Sc}{\mathsf S}
\newcommand{\Rc}{\mathsf{rec}}
\newcommand{\Oc}{\mathsf{oracle}}
\newcommand{\Boost}{\mathsf{Boost}}
\newcommand{\Step}{\mathsf{Step}}
\newcommand{\Entry}{\mathfrak N}
\newcommand{\Cfive}{C_5}
\newcommand{\code}[1]{\texttt{#1}}
\newcommand{\sem}[1]{\llbracket #1\rrbracket}
\lstset{basicstyle=\ttfamily\small,columns=fullflexible,keepspaces=true,breaklines=true,
 frame=single,framerule=.3pt,rulecolor=\color{black!25},showstringspaces=false,
 xleftmargin=5pt,xrightmargin=5pt,aboveskip=8pt,belowskip=8pt}

\title{A Constructive Lean Contender\texorpdfstring{\\}{ }Beyond Contender 5\\[7pt]
\large Ramified W-reflection, higher-type feedback,\\and an explicit tower-height separation}
\author{Research prepared for Vladimir Reshetnikov\\\small OpenAI}
\date{12 October 2026}

\begin{document}
\maketitle
\begin{abstract}
We propose a natural number $\Entry=R_\Theta(64)$ for a Lean adaptation of the
\emph{name-the-biggest-number} contest. Its definition combines a finite
maximum over a total language, recursion on countably branching well-founded
trees, diagonalization through arbitrary finite function types, and feedback
from previously defined numerical values into type-level indices. The
construction does not refer to the preceding winning constant and does not use
a halting oracle, unrestricted semantic maximization, or an assumed universal
Lean evaluator.

We give an exact semantic port of the upstream Contender 5, including its
behavior on malformed raw terms. A short quotation argument proves
$\Cfive<\Entry$. A stronger argument, using a depth-20 oracle-iteration term,
proves $E_{\Cfive+1}(64)<\Entry$, where $E_0(x)=x$ and
$E_{j+1}(x)=2^{E_j(x)}$. We also establish explicit child-stage comparisons,
a $2^{28}$-branch visibility bound at input 64, and eventual domination of
every fixed numerical program in a stage's own language.

The accompanying artifact contains Lean definition and proof source, a
strict axiom/timing gate, and a Python small-instance model with 62,065 passing
assertions. \textbf{The Lean source has not been compiler-checked: the
preparation environment had no Lean toolchain and attempted downloads failed.}
Consequently, kernel acceptance, the empty-axiom audit, and the contest's
15-second/60-second timing requirements remain unverified. The mathematical
proposal and proofs are not presented as an already accepted contest entry.
\end{abstract}

\clearpage
\setcounter{tocdepth}{1}
\tableofcontents
\clearpage

\section{What is being proposed, and what is being claimed}

The user asks for a serious constructive contender after replacing Coq with
Lean in Cody Roux's contest. The source repository attributes Contender 5 to
Samuel Gr\"utter and defines it as the largest numerical output at weighted
syntax depth 42 in a totalized simply typed language with natural-number
recursion~\cite{repo,contender}. The original rules concern \emph{typechecking
time}, not the time required to evaluate the resulting natural number.
They also prohibit axioms and defining the next entry from the previous
winning constant~\cite{rules}.

Our proposed declaration is
\begin{equation}\label{eq:entry}
 \boxed{\Entry=R_\Theta(64).}
\end{equation}
Every symbol in this equation will receive a constructive definition. In
particular, $R$ is a tree recursor, not an appeal to a vaguely specified
fast-growing hierarchy. The central numerical comparison is
\begin{equation}\label{eq:headline}
 \boxed{E_{\Cfive+1}(64)<\Entry,\qquad
 E_0(x)=x,\quad E_{j+1}(x)=2^{E_j(x)}.}
\end{equation}
The role of $\Cfive$ in \eqref{eq:headline} is \emph{only in the theorem}; it
is absent from the definition of $\Entry$.

There is no claim of a mathematically greatest constructive natural number.
Any fixed total definition admits further amplification. Nor do we assert a
sharp proof-theoretic ordinal for this construction. The contribution is a
particular ambitious, fully specified proposal with provable separation,
carefully delimited formalization status, and reproducible artifacts.
The general finite-language-maximization idea is inherited from the upstream
entry. Our proposed combination of W-recursive object-language operations,
ramified numerical oracles, type-index feedback, and explicit budgeted
comparison witnesses is developed here; a comprehensive priority claim is
not made.

\subsection{A precise Lean adaptation of the rules}

We retain the mathematical intent of the original contest and propose the
following operational interpretation. The pinned target is Lean 4.35.0;
pinning is necessary for reproducibility, not a claim that this source was
successfully checked on that version~\cite{leanrelease}.

\begin{enumerate}[label=\arabic*.]
\item The entry is a closed definition of type \code{Nat}, using only total
computational definitions. It must not depend on a custom axiom, a proof
placeholder, classical choice, noncomputable selection, an unsafe replacement,
or an external numerical oracle.
\item The new definition dependency closure must compile in at most 15
seconds, and the comparison proof in at most 60 seconds, on a disclosed
reasonable machine. Distributed standard-library artifacts may be prebuilt;
new contestant modules may not be hidden in a precompiled cache.
\item The entry's definition must not refer to the previous winning
\emph{number}. Reusing an explicitly charged source language and its evaluator
is allowed. The comparison proof may name the previous entry.
\item The transitive axiom reports for the entry and required comparison
must be empty. Importing a library that contains unused axioms is not itself
an axiom dependency; using one is.
\end{enumerate}

Lean distinguishes computational definitions from additional principles such
as classical choice and propositional extensionality. These distinctions make
an actual transitive audit important~\cite{leanaxioms}. The package's gate
intentionally does not whitelist such principles. The surrounding explanatory
mathematics occasionally uses extensional descriptions of functions; those
analysis-only identities are not imported into the strict Lean proof path.

\subsection{Validation boundary}

Three different statements must not be conflated: a complete mathematical
proof, proposed formal proof source, and a successful kernel run. This report
provides the first two for the basic strict comparison. It does not provide
the third. The stronger tower theorem, the binary budget bounds, and the
higher-type extensional identities are proved in this article and tested in
small instances where appropriate; they are not claimed as completed Lean
formalizations. Section~\ref{sec:validation} records the exact evidence.

\section{The upstream number, without changing its semantics}
\label{sec:legacy}

\subsection{Types, raw terms, and weighted depth}

Let the old types be generated by
\[
 A,B ::= \tyN \mid A\to B,
 \qquad \sem{\tyN}=\N,\qquad \sem{A\to B}=\sem A\to\sem B.
\]
Their depth is $\td(\tyN)=1$ and
$\td(A\to B)=1+\max(\td A,\td B)$.
The raw terms are
\[
 t ::= \mathsf{var}(i)\mid\lambda_{A,B}t\mid t\,u
       \mid\mathsf 0\mid\Sc\mid\Rc_A .
\]
The recursor has semantic type
\[
 \Rc_A:\sem A\to(\N\to\sem A\to\sem A)\to\N\to\sem A.
\]
The weighted term depth is
\begin{align*}
 \depth(\mathsf{var}(i))&=i+2,\\
 \depth(\lambda_{A,B}t)&=1+\max(\td A,\td B,\depth t),\\
 \depth(tu)&=1+\max(\depth t,\depth u),\\
 \depth(\mathsf0)=\depth(\Sc)&=1,\qquad
 \depth(\Rc_A)=1+\td A.
\end{align*}
In particular, depth counts type annotations and variable labels. It is not
ordinary node count or source-byte length.

\subsection{The significance of malformed terms}

The upstream evaluator does \emph{not} first discard terms that fail a typing
judgment. It returns a tagged semantic value for every raw term. This is
visible directly in \code{interp\_term}, \code{cast\_impl}, and
\code{error} in the source~\cite{contender}. Replacing this evaluator with a
well-typed-term enumerator would silently change the target of comparison.

For each type define its default value by
\[
 \bot_{\tyN}=0,\qquad \bot_{A\to B}(x)=\bot_B.
\]
A computational cast pair between two types is identity at matching natural
types. At matching arrow shapes with syntactically equal component types it
uses recursively defined contravariant input and covariant output maps. If
these tests fail, both maps return the appropriate default value. Thus no
opaque equality proof is needed to run a cast.

An environment is a list of tagged values. Variable labels are absolute
binder levels: label 0 refers to the outermost binder, not the newest one.
For an environment of length $l$, lookup of $i<l$ uses position $l-1-i$ in the
cons-based list; otherwise lookup fails and produces the tagged value
$(\tyN,0)$. A lambda creates a semantic function, evaluates its body under a
new tagged argument, and casts the body's result to its stated codomain.
Application of a tagged function casts its argument to the function's domain.
Application of a nonfunction produces $(\tyN,0)$.

Finally, $\val_0(t)$ evaluates $t$ in the empty environment and returns its
natural component if its tag is $\tyN$; otherwise it returns zero. This
includes raw terms with free variable labels: ``empty-environment
interpretation'' is not a promise of syntactic closedness. For example,
$\val_0(\Sc\,\Sc)=1$, since the function-valued argument is cast to natural
zero rather than rejected.

\subsection{The finite enumerator and the selected maximum}

Let $U_0=[]$ and
\[
 U_{d+1}=[\tyN]\mathbin{+\!+}[A\to B\mid A,B\in U_d].
\]
Let $T_0=[]$. The list $T_{d+1}$ concatenates, in the upstream order,
\begin{align*}
 &[\mathsf{var}(i)\mid i=d-1,d-2,\ldots,0],\\
 &[\lambda_{A,B}t\mid A,B\in U_d,\ t\in T_d],\\
 &[tu\mid t,u\in T_d],\quad [\mathsf0],\quad[\Sc],\quad
 [\Rc_A\mid A\in U_d].
\end{align*}
The variable list is empty for $d=0$. Induction establishes that membership in
$T_d$ is equivalent to having depth at most $d$; the upstream source proves
this exact characterization~\cite{contender}.

For fidelity, the Lean baseline uses the original argmax algorithm, not just
an independently specified maximum. Start with current value 0 and current
best term $\mathsf0$, scan the list, and replace the best term exactly when a
strictly larger value appears. Write the resulting selected term as
\[
 t_* = \mathsf{maxBy}(\val_0,0,\mathsf0,T_{42}),\qquad
 \Cfive=\val_0(t_*).
\]

\begin{lemma}[Selection and maximum]\label{lem:selection}
The term $t_*$ belongs to $T_{42}$, and
$\Cfive=\max(\{0\}\cup\{\val_0(t):t\in T_{42}\})$.
\end{lemma}
\begin{proof}
At each scan step the selected term is either the old seed or a term already
visited. The seed $\mathsf0$ belongs to $T_{42}$, so the final selection is in
that list. The invariant ``current value equals the selected value and is the
maximum of the seed value and all visited values'' holds initially and is
preserved by either branch of the strict comparison. It yields the displayed
identity after the final element.
\end{proof}

\begin{proposition}[Mathematical fidelity of the port]\label{prop:port}
The computational definitions in \code{Legacy.lean}, together with
\code{Baseline.lean}, denote the same natural number as the pinned upstream
\code{contender\_5}.
\end{proposition}
\begin{proof}
Match type constructors, raw constructors, natural numbers, and lists
componentwise. Relate semantic natural values by equality and arrow values
by preservation of this relation on related inputs. Structural induction on
types shows that defaults are related and that both components of each cast
pair preserve the relation. The same Boolean comparisons select the same
cast branches. The variable lookup functions use identical lengths,
subtraction, and zero-based list positions.

Now induct on the raw term, with the induction hypothesis quantified over
related environments. The lambda case extends those environments by related
arguments and uses the codomain cast lemma. The application case first uses
the two term hypotheses, then the domain cast lemma. Natural recursion
preserves the relation by induction on its natural argument. Constants and
failed lookups are immediate. Thus empty-environment numerical observations
agree for every raw term. The two list enumerators agree by induction on the
depth. Finally, the argmax scans use the same strict comparison on equal
natural values, start from corresponding seeds, and make corresponding
choices. Their numerical results therefore agree.
\end{proof}

This is a semantic port argument, not a claim of a machine-checked
Coq-to-Lean translation certificate. The upstream Coq project was not compiled
in this environment. The checked-in provenance pins both its commit and the
exact source blob.

\section{A stronger, but still total, finite language}
\label{sec:language}

\subsection{The tree type}

Define the strictly positive inductive type
\begin{equation}\label{eq:tree}
 \W ::= \leaf \mid \Sup(\N\to\W).
\end{equation}
An element is a countably branching well-founded tree. Its branch function
may be represented by a finite program; it is not an explicitly stored
infinite list. The constructor $\Sup$ is \emph{not} an ordinal supremum and we
impose no quotient identifying trees of equal ordinal rank.

For any result type $X$, its recursor satisfies
\begin{align}
 \mathsf{fold}_X(z,s,\leaf)&=z,\label{eq:wfold}\\
 \mathsf{fold}_X(z,s,\Sup f)&=
 s\bigl(f,\ i\mapsto\mathsf{fold}_X(z,s,f(i))\bigr).
\end{align}
The recursive values are supplied as a function of the child index, not
computed by eager traversal of every child. Strict positivity and recursors
are the relevant Lean mechanisms~\cite{leaninductive,leanrecursive}.

\subsection{Types and primitives}

The new type language is
\[
 A,B ::= \tyN\mid\tyW\mid A\to B,
 \qquad\sem{\tyW}=\W.
\]
There is a natural recursor and a tree recursor at \emph{every} result type
in this finite grammar. For a fixed binary function $q:\N\to\N\to\N$, the
additional primitives are
\begin{align*}
 \mathsf{treeZero}&:\W,&
 \mathsf{treeSup}&:(\N\to\W)\to\W,\\
 \Oc&:\N\to\N\to\N,&\sem{\Oc}&=q,\\
 \mathsf{treeRec}_A&:\sem A\to
 ((\N\to\W)\to(\N\to\sem A)\to\sem A)\to\W\to\sem A.
\end{align*}
The ordinary zero, successor, variables, lambdas, applications, and natural
recursors are retained. Defaults and computational casts are extended to
$\tyW$ by $\bot_{\tyW}=\leaf$ and identity on matching tree types.
Nonfunction application again gives tagged natural zero.

There is also a constructor $\Qte(t)$ for \emph{any raw old term} $t$.
It evaluates $t$ in the old empty environment and returns the tagged natural
$\val_0(t)$. Its depth charge is $1+\depth(t)$. This is an explicit embedding
of the old numerical semantics, not a constant for the old winning number,
not a call to the old maximum, and not a free abbreviation for an arbitrarily
large source file.

Let $\val_q(e)$ denote the new empty-environment natural observation. Every
raw new term has an observation; no undecidable validity filter is inserted.
The only oracle during this section is a \emph{parameter}. In the eventual
entry every such parameter is replaced by an already structurally smaller,
total function.

\subsection{Depth and exact finite enumeration}

Both ground types have type depth 1; arrow depth is unchanged. Ordinary term
depths are unchanged, all five unannotated constants
$\mathsf0,\Sc,\mathsf{treeZero},\mathsf{treeSup},\Oc$ have depth 1, each
recursor annotated by $A$ has depth $1+\td A$, and
\[
 \depth(\Qte(t))=1+\depth(t).
\]
The new type list is
\[
 V_0=[],\qquad V_{d+1}=[\tyN,\tyW]
       \mathbin{+\!+}[A\to B\mid A,B\in V_d].
\]
The new term list $\mathcal T_0=[]$ and $\mathcal T_{d+1}$ is the concatenation
of the following eight blocks:
\begin{equation}\label{eq:enum}
\begin{split}
 &\mathcal T_d;\quad
 [\mathsf0,\Sc,\mathsf{treeZero},\mathsf{treeSup},\Oc];\quad
 [\mathsf{var}(i):i<d];\\
 &[\lambda_{A,B}e:A,B\in V_d,\ e\in\mathcal T_d];\quad
 [ef:e,f\in\mathcal T_d];\\
 &[\Rc_A:A\in V_d];\quad
 [\mathsf{treeRec}_A:A\in V_d];\quad
 [\Qte(t):t\in T_d].
\end{split}
\end{equation}
The order of variable labels in the implementation is descending; the other
products use list order. The first block deliberately makes the enumerator
cumulative. Repetition changes execution cost but never changes a maximum.

\begin{lemma}[Coverage]\label{lem:coverage}
The list $\mathcal T_d$ is finite and its elements are exactly the raw terms
of depth at most $d$. Moreover, $\mathcal T_d$ is included, with occurrences,
in $\mathcal T_{d+1}$.
\end{lemma}
\begin{proof}
Finiteness follows by induction, since each step uses finite products,
maps, and concatenations. For types, induction gives
$A\in V_d\iff\td A\le d$. For terms at step $d+1$, inspect the outer
constructor. An application has both children in $\mathcal T_d$ exactly when
its depth is at most $d+1$; a lambda additionally needs both annotations in
$V_d$. The analogous equivalences handle recursors, variable labels, and
quotation via the old enumerator. Constants occur at every positive step.
The carried-over block cannot introduce a term of excessive depth. These
observations prove both directions of coverage inductively. The last claim
is immediate from the first block.
\end{proof}

The strict Lean comparison uses direct occurrence lemmas for the particular
constructors it needs, so it does not depend on a formalized converse to this
coverage lemma.

\begin{definition}[Finite semantic envelope]\label{def:envelope}
For a total binary $q$, define
\begin{equation}\label{eq:envelope}
 M_q(d)=\max\bigl(\{0\}\cup
       \{\val_q(e):e\in\mathcal T_d\}\bigr).
\end{equation}
The implementation uses a recursive list maximum, not classical selection.
\end{definition}

\begin{proposition}[Totality and monotonicity]
For every total computable $q$, $M_q$ is total computable, and
$M_q(d)\le M_q(d+1)$. Every occurrence $e\in\mathcal T_d$ satisfies
$\val_q(e)\le M_q(d)$.
\end{proposition}
\begin{proof}
The interpreter recurses on raw syntax. Semantic recursion inside its
primitive constants is natural recursion or the tree recursor
\eqref{eq:wfold}. Both have total input functions of the indicated types.
Thus every list element has a natural observation. A finite maximum of these
observations is computed by list recursion. Occurrence bounds follow by
induction on the maximum computation, and monotonicity follows from the
cumulative block in \eqref{eq:enum}.
\end{proof}

The theorem is a totality statement, not an efficiency statement. In
particular, evaluating $M_q(68)$ is not part of the proposed validation run.

\section{Short witnesses and escaping a fixed language}
\label{sec:diagonal}

\subsection{Unary witnesses}

Write $\overline n=\Sc^n\mathsf0$ for the unary numeral. Then
$\depth(\overline n)=n+1$ and $\val_q(\overline n)=n$.
For $i\le n$, the raw term
\begin{equation}\label{eq:unarywitness}
 w_{i,n}=\Sc\bigl(\Oc\,\overline i\,\overline n\bigr)
\end{equation}
has depth
$\max(i+4,n+3)\le n+4$ and value $q(i,n)+1$. Therefore
\begin{equation}\label{eq:oraclebasic}
 q(i,n)<M_q(n+4)\quad(i\le n).
\end{equation}
Likewise, the numeral $\overline{n+1}$ proves
\begin{equation}\label{eq:inputgrowth}
 n<M_q(n+4).
\end{equation}
These statements need no monotonicity assumption on $q$.

\subsection{Binary numerals with an explicit syntax cost}

Because variable names are absolute binder levels, embedding a term below
new binders requires care. At top level define the depth-8 term
\begin{equation}\label{eq:double}
 D=\Rc_{\tyN}\,\mathsf0\,
       (\lambda_{\tyN,\tyN\to\tyN}
         \lambda_{\tyN,\tyN}\Sc(\Sc(\mathsf{var}(1)))).
\end{equation}
It denotes the function $x\mapsto2x$. The outer step argument has level 0;
the accumulator has level 1. The natural-recursion equations prove the
claim by induction on the input.

Let $b(0)=\mathsf0$. For $n>0$, start with $\mathsf0$ and read the ordinary
binary digits of $n$ from left to right, replacing the current term $e$ by
$De$ for digit 0 and by $\Sc(De)$ for digit 1. These occurrences of $D$ are
all in the empty environment. With $\bits(0)=0$ and
$\bits(n)=\lfloor\log_2 n\rfloor+1$ for $n>0$, we have
\begin{equation}\label{eq:binary}
 \val_q(b(n))=n,\qquad \depth(b(n))\le8+2\bits(n).
\end{equation}
The value claim follows from binary positional notation. The first digit
creates depth at most 10; every later digit increases depth by at most two.
The displayed bound also covers zero.

\begin{theorem}[Explicit escape from a fixed oracle language]
\label{thm:escape}
Fix $q$. Let $p$ be any fixed raw term whose empty-environment interpretation
has type $\tyN\to\tyN$ and denotes a function $f:\N\to\N$. Put
$c=\depth(p)$. Then
\begin{equation}\label{eq:escape}
 f(n)<M_q(n)\qquad\text{for every }n\ge\max(c+2,32).
\end{equation}
Consequently, $M_q$ is not denoted by any such fixed unary term in its own
$q$-language.
\end{theorem}
\begin{proof}
The term $\Sc(p\,b(n))$ has value $f(n)+1$ and depth at most
$2+\max(c,8+2\bits(n))$. For $n\ge32$,
\[
 10+2\bits(n)\le n.
\]
Indeed, on the block $2^r\le n<2^{r+1}$ with $r\ge5$, its left side is
$12+2r\le2^r\le n$. The inequality at $r=5$ holds, and its preservation
when $r$ increases is immediate. The assumption $n\ge c+2$ handles the
other term in the maximum. The witness therefore belongs to
$\mathcal T_n$, proving \eqref{eq:escape}. If a fixed $p$ denoted $M_q$,
\eqref{eq:escape} would assert $M_q(n)<M_q(n)$ for sufficiently large $n$.
\end{proof}

The conclusion concerns the specified finite object language, not all Lean
programs. The host can define $M_q$ because it can interpret the whole type
code family and recurse on syntax. An object-language term has one fixed
finite type. No internal global reflection principle is inferred.

\subsection{Binary visibility at a fixed budget}

Replacing both unary numerals in \eqref{eq:unarywitness} by binary numerals
gives depth at most
\begin{equation}\label{eq:binarybudget}
 \max\bigl(11+2\bits(i),\ 10+2\bits(n)\bigr).
\end{equation}
This follows by counting the two curried applications and the outer
successor: the exact cost is
$\max(\depth(b(i))+3,\depth(b(n))+2)$.
At $n=64$, every $i<2^{28}$ fits within depth 68, since the two upper bounds
are at most 67 and 24. Thus
\begin{equation}\label{eq:binaryvisible}
 q(i,64)<M_q(68)\qquad(0\le i<2^{28}).
\end{equation}
This is a sufficient bound, not an optimal coding claim. Some larger labels
also fit because their binary expansions contain fewer 1s.

\section{Ramified reflection on trees}
\label{sec:reflection}

\begin{definition}[Stage functions]
Define $R:\W\to\N\to\N$ by the tree recursor:
\begin{align}
 R_{\leaf}(n)&=M_{q_0}(n+4),&q_0(i,x)&=0,\label{eq:reflect0}\\
 R_{\Sup f}(n)&=M_{q_f}(n+4),&q_f(i,x)&=R_{f(i)}(x).
 \label{eq:reflectsup}
\end{align}
\end{definition}
The recursive clause is exactly the step algebra for \eqref{eq:wfold} with
result type $\N\to\N$. The tree children need not have a uniform numerical
height bound. Their structural accessibility is supplied by the inductive
type itself.

\begin{theorem}[Constructive totality]\label{thm:totalR}
Every $R_a$ is total. For any closed computably constructed $a:\W$, the
function $R_a:\N\to\N$ is computable in principle.
\end{theorem}
\begin{proof}
The leaf algebra in \eqref{eq:reflect0} is total. Given the recursive
functions $i\mapsto R_{f(i)}$, the step algebra merely supplies them as the
binary oracle to a total finite envelope. Hence it is a total function
$\N\to\N$. Applying the tree recursor gives the result. All operations
involved are computational recursors or structural definitions, so a closed
constructive input yields a closed computational function.
\end{proof}

The phrase ``computably constructed'' matters when interpreting the theorem
in ordinary set-theoretic semantics: arbitrary mathematical functions
$\N\to\W$ can carry noncomputable information. Our entry uses only the
explicit closed constructions below; it does not select arbitrary semantic
trees from a full set-theoretic function space.

\begin{theorem}[Visible-child separation]\label{thm:child}
For every $f:\N\to\W$ and every $i\le n$,
\[
 R_{f(i)}(n)<R_{\Sup f}(n).
\]
Also $n<R_a(n)$ for every tree $a$ and natural $n$.
\end{theorem}
\begin{proof}
Apply \eqref{eq:oraclebasic} to $q_f$ and use
\eqref{eq:reflectsup}. For the second statement, apply
\eqref{eq:inputgrowth} to the oracle at $a$, treating the leaf and supremum
constructors separately.
\end{proof}

\begin{corollary}[Many visible branches at input 64]\label{cor:many}
For every $i<2^{28}$,
$R_{f(i)}(64)<R_{\Sup f}(64)$.
\end{corollary}
\begin{proof}
Use \eqref{eq:binaryvisible} with $q=q_f$.
\end{proof}

At each fixed tree $a$, Theorem~\ref{thm:escape} shows that $R_a$ eventually
outgrows every fixed unary program in the language using the local oracle
assigned to $a$. This applies at the leaf as well as at nonleaf stages. It
is a uniform syntactic phenomenon, not an assertion derived only from a
suggested ordinal rank.

\section{A quantitative amplification: oracle-fed iteration}
\label{sec:towerengine}

A single reflection step can do substantially more than add one to a child
value. We now prove the numerical amplification that yields
\eqref{eq:headline}.

\subsection{A depth-14 power-of-two program}

Let $D^{[l]}$ be the doubling program \eqref{eq:double} with its accumulator
variable changed from level 1 to level $l+1$, so it denotes doubling when
evaluated in an environment of length $l$. Its depth is $l+8$. Define the
closed program
\begin{equation}\label{eq:powprogram}
 P=\Rc_{\tyN}\,(\Sc\mathsf0)\,
 \bigl(\lambda_{\tyN,\tyN\to\tyN}
       \lambda_{\tyN,\tyN}
       (D^{[2]}\,\mathsf{var}(1))\bigr).
\end{equation}
It has depth 14 and denotes $x\mapsto2^x$. The change to level 3 inside
$D^{[2]}$ is essential: simply inserting the original $D$ below two new
binders would capture the wrong accumulator under the upstream convention.

To verify the depth, $\depth(D^{[2]})=10$, application to
$\mathsf{var}(1)$ gives 11, the two outer lambdas give 13, and the final
application of the partially applied recursor gives 14. Its value follows
from the natural-recursion equations: start at one and double the
accumulator at each step.

\begin{lemma}[Uniform exponential lower bound]\label{lem:exponential}
For every tree $a$ and every $x\ge32$,
\[
 2^x\le R_a(x).
\]
\end{lemma}
\begin{proof}
The program $P\,b(x)$ is independent of the local oracle and has value $2^x$.
Its depth is at most $1+\max(14,8+2\bits(x))\le x+4$ for $x\ge32$, by the
binary inequality proved in Section~\ref{sec:diagonal}. It therefore occurs
in the envelope defining $R_a(x)$.
\end{proof}

\subsection{The feedback witness}

For a fixed binary $q$, write $g(x)=q(0,x)$. In the empty environment, let
\[
 A_q=\lambda_{\tyN,\tyN\to\tyN}\lambda_{\tyN,\tyN}
      (\Oc\,\mathsf0\,\mathsf{var}(1)).
\]
The syntax $A_q$ uses the oracle constant, not the semantic function $q$ as
an extra constructor. Its depth is 6, and it ignores its first argument
while applying $g$ to its accumulator. Define
\begin{equation}\label{eq:feedbackwitness}
 F_n=\Sc\Bigl(
    \Rc_{\tyN}\,b(n)\,A_q\,
         (\Oc\,\mathsf0\,b(n))\Bigr).
\end{equation}
Writing $g^{[j]}$ for $j$-fold iteration, the recursor equations give
\begin{equation}\label{eq:feedbackvalue}
 \val_q(F_n)=g^{[g(n)]}(n)+1.
\end{equation}
The two occurrences of $b(n)$ remain at top level; the only variable within
the step function is the correctly scoped accumulator at level 1.

\begin{lemma}[Feedback witness cost]\label{lem:feedbackcost}
The term $F_n$ has depth at most
$\max(9,12+2\bits(n))$. For $n\ge32$ this is at most $n+4$.
For the specified binary encoding, $F_{64}$ has exact depth 20.
\end{lemma}
\begin{proof}
Put $d=\depth(b(n))$. The recursor applied to its initial value has depth
$1+\max(2,d)$. Applying it to the depth-6 step gives depth
$\max(7,d+2)$. The iteration-count term
$\Oc\,\mathsf0\,b(n)$ has depth $\max(3,d+1)$. The next application and
outer successor give total depth $\max(9,d+4)$. Insert
\eqref{eq:binary}. For $n\ge32$, the blockwise estimate used earlier gives
the required bound. The binary representation of 64 is $1000000$: its
numeral has depth 16, so $\max(9,16+4)=20$.
\end{proof}

\begin{theorem}[Oracle-fed iteration separation]\label{thm:feedback}
For every $f:\N\to\W$, set $g=R_{f(0)}$. Then for every $n\ge32$,
\begin{equation}\label{eq:feedbackseparation}
 g^{[g(n)]}(n)<R_{\Sup f}(n).
\end{equation}
\end{theorem}
\begin{proof}
Use $q(i,x)=R_{f(i)}(x)$ in \eqref{eq:feedbackwitness}. The witness has
value one more than the left side by \eqref{eq:feedbackvalue}, and fits the
stage's $n+4$ budget by Lemma~\ref{lem:feedbackcost}. The envelope bounds
its value.
\end{proof}

\begin{corollary}[Tower separation from the zeroth child]\label{cor:tower}
For $n\ge32$ and $m=R_{f(0)}(n)$,
\[
 E_m(n)<R_{\Sup f}(n).
\]
\end{corollary}
\begin{proof}
Lemma~\ref{lem:exponential} shows $g(x)\ge2^x$ on $x\ge32$. Induction on
$j$ gives $g^{[j]}(n)\ge E_j(n)$: the base is equality, and at a successor
the induction hypothesis ensures the input remains at least 32, after which
$g(g^{[j]}(n))\ge2^{g^{[j]}(n)}\ge2^{E_j(n)}$. Set $j=m$ and apply
Theorem~\ref{thm:feedback}.
\end{proof}

These proofs use explicit finite terms rather than guesses about growth-rate
classifications. The critical cost is the small \emph{description} of the
iteration count $g(n)$, not the size of that count.

\section{Higher-type countable iteration}
\label{sec:higher}

We next construct the tree index $\Theta$. This adds substantial structure
beyond what is needed for the basic comparison theorem.

\subsection{Finite levels and pointwise suprema}

Define a family of types
\begin{equation}\label{eq:levels}
 L_0=\W,\qquad L_{k+1}=L_k\to L_k.
\end{equation}
Define a countable-supremum operation at each level by
\begin{align}
 S_0(f)&=\Sup(f),\label{eq:suplevel}\\
 S_{k+1}(f)(x)&=S_k(j\mapsto f(j)(x)).
\end{align}
For $u:L_k\to L_k$, let $u^{[0]}(x)=x$ and
$u^{[j+1]}(x)=u(u^{[j]}(x))$. The countable-iteration operator is
\begin{equation}\label{eq:omega}
 \Omega_k(u)(x)=S_k(j\mapsto u^{[j]}(x)),
 \qquad \Omega_k:L_{k+2}.
\end{equation}
The family $L_k$ is a natural-number-indexed family of types. It is not a
sequence of stronger logical axioms. Every member lies in the same small
universe, and Lean can express the family by natural recursion.

For a starting tree $a$ and an operation $s:\W\to\W$, set
\begin{equation}\label{eq:seed}
 \sigma_0=a,\qquad\sigma_1=s,\qquad
 \sigma_{k+2}=\Omega_k.
\end{equation}
These seeds have types $\sigma_k:L_k$. Define evaluation down the levels by
\begin{equation}\label{eq:collapse}
 C_0(x)=x,\qquad C_{k+1}(u)=C_k(u(\sigma_k)),
\end{equation}
with dependence on $a,s$ understood. Finally let
\begin{equation}\label{eq:high}
 H_k(a;s)=C_{k+2}(\Omega_k).
\end{equation}
Every operation here is a structural recursion on a natural number, a finite
function iteration, or a tree constructor.

\subsection{Structural identities and their meaning}

\begin{lemma}[Commutation with countable suprema]\label{lem:commute}
In extensional mathematical notation, for every $f:\N\to L_k$,
\[
 C_k(S_k(f))=\Sup(j\mapsto C_k(f(j))).
\]
\end{lemma}
\begin{proof}
For $k=0$ the two sides are identical. At level $k+1$, apply
\eqref{eq:collapse} and then \eqref{eq:suplevel} to obtain
$C_k(S_k(j\mapsto f(j)(\sigma_k)))$. The induction hypothesis identifies
this with $\Sup(j\mapsto C_k(f(j)(\sigma_k)))$, which is the claimed
expression by the definition of $C_{k+1}$.
\end{proof}

\begin{proposition}[Visible predecessor in adjacent levels]\label{prop:adjacent}
One has
\begin{align}
 H_0(a;s)&=\Sup(j\mapsto s^{[j]}(a)),\label{eq:high0}\\
 H_k(a;s)&=\Sup(j\mapsto C_k(\sigma_{k+1}^{[j]}(\sigma_k))).
 \label{eq:highroot}
\end{align}
In the root presentation of $H_{k+1}(a;s)$ in \eqref{eq:highroot}, child 1 is
$H_k(a;s)$. Consequently, for $n\ge1$,
\[
 R_{H_k(a;s)}(n)<R_{H_{k+1}(a;s)}(n).
\]
\end{proposition}
\begin{proof}
Expanding two steps of collapse in \eqref{eq:high} gives
$C_k(\Omega_k(\sigma_{k+1})(\sigma_k))$. Expand \eqref{eq:omega} and apply
Lemma~\ref{lem:commute} to obtain \eqref{eq:highroot}. Its $k=0$ instance
is \eqref{eq:high0}. At level $k+1$, the child at index 1 is
\[
 C_{k+1}(\sigma_{k+2}(\sigma_{k+1}))
 =C_{k+1}(\Omega_k(\sigma_{k+1}))=H_k(a;s).
\]
Theorem~\ref{thm:child} now applies with child index 1.
\end{proof}

The general equality statements in this subsection are analysis-only
extensional identities. The Lean proof source includes \eqref{eq:high0},
which is definitional, but does not claim a completed empty-axiom
formalization of Lemma~\ref{lem:commute} or the general adjacent-level
identity. A pointwise logical relation is one possible way to formalize
these statements without adding functional extensionality as a dependency.
None of these analysis-only equalities is needed to prove the entry's
required strict comparison.

\section{Numerical feedback into type-level indices}
\label{sec:grand}

Define a simple tree successor operation
\[
 \Step(a)=\Sup(i\mapsto a).
\]
Unlike a quotient ordinal successor, this is merely an explicit tree with
the same child at every index. The name serves as a structural mnemonic.

\begin{definition}[Feedback boost]\label{def:boost}
Define
\begin{equation}\label{eq:boost}
 \Boost(a)=\Sup(b_a),\qquad
 b_a(0)=a,\quad
 b_a(k+1)=H_{R_a(k+64)}(a;\Step).
\end{equation}
\end{definition}
The construction feeds a reflected \emph{number} into the height of the
finite-type family \eqref{eq:levels}. This is stronger structurally than
choosing a larger fixed literal in a single envelope. It does not rely on
an assumption that this type height has a particular ordinal name.

There is no circular definition. The operation $R$ was completed in
Section~\ref{sec:reflection} before $\Boost$ was introduced. The body of
$\Boost$ calls that already defined operation on its input $a$. The
functions $H_k$ were likewise defined for all natural $k$ before any
$R_a(k+64)$ was used as their index.

\begin{theorem}[Feedback comparisons]\label{thm:boost}
For every $a$ and every $n$,
\[
 R_a(n)<R_{\Boost(a)}(n).
\]
For $n\ge1$ one also has
\[
 R_{H_{R_a(64)}(a;\Step)}(n)<R_{\Boost(a)}(n).
\]
\end{theorem}
\begin{proof}
The retained branch at index 0 gives the first inequality via
Theorem~\ref{thm:child}. Branch 1 is exactly
$H_{R_a(64)}(a;\Step)$, giving the second. No comparison between the
numerical oracles as pointwise functions is required.
\end{proof}

\begin{definition}[The grand index and proposed entry]
Set
\begin{equation}\label{eq:grand}
 \Theta=\Sup(k\mapsto H_k(\leaf;\Boost)),\qquad
 \Entry=R_\Theta(64).
\end{equation}
\end{definition}

\begin{theorem}[Some explicit stages exceeded by the entry]\label{thm:grand}
For $k\le n$,
\[
 R_{H_k(\leaf;\Boost)}(n)<R_\Theta(n).
\]
For $j\le n$,
\[
 R_{\Boost^{[j]}(\leaf)}(n)<R_\Theta(n).
\]
At $n=64$, both conclusions hold for every $k<2^{28}$ and every $j<2^{28}$,
respectively.
\end{theorem}
\begin{proof}
The first statement uses child $k$ of \eqref{eq:grand}. The second uses
child $j$ of $H_0(\leaf;\Boost)$ in \eqref{eq:high0}, followed by child 0
of $\Theta$, and composes the two strict inequalities. For input 64, use
Corollary~\ref{cor:many} instead of the unary-index version at the indicated
child selections.
\end{proof}

The word ``all'' in ``all finite levels'' describes the index construction:
$\Theta$ contains a child for each natural level. A fixed numerical input
has a fixed syntax budget. The theorem carefully states which branch
comparisons that budget proves, rather than claiming that one fixed output
exceeds an unbounded sequence of natural outputs at that input.

\section{The exact comparison with Contender 5}
\label{sec:comparison}

\begin{theorem}[Uniform comparison with the upstream baseline]
\label{thm:baseline}
For every total binary oracle $q$,
\[
 \Cfive<M_q(68).
\]
Consequently, for every tree $a$,
\[
 \Cfive<R_a(64).
\]
\end{theorem}
\begin{proof}
By Lemma~\ref{lem:selection}, $t_*\in T_{42}$. The quotation block of
\eqref{eq:enum} places $\Qte(t_*)$ in $\mathcal T_{43}$. Both this term
and $\Sc$ therefore occur at depth 43, so the application
\[
 e_* = \Sc(\Qte(t_*))
\]
occurs in $\mathcal T_{44}$. Cumulativity puts it in $\mathcal T_{68}$.
Its value is $\val_0(t_*)+1=\Cfive+1$, independently of $q$. Thus
$\Cfive+1\le M_q(68)$. For any tree $a$, its stage at 64 is exactly such
an envelope, since $64+4=68$; distinguish the two constructors in
\eqref{eq:reflect0}--\eqref{eq:reflectsup}.
\end{proof}

\begin{corollary}[Required strict inequality]\label{cor:entry}
The proposed number is strictly greater than the pinned upstream Contender 5:
$\Cfive<\Entry$.
\end{corollary}
\begin{proof}
Take $a=\Theta$ in Theorem~\ref{thm:baseline}.
\end{proof}

This proof does not evaluate the old maximum or the new one. It proves that
a particular \emph{description} occurs in the relevant finite list. The
successor operation appears in that proof witness, not as the definition of
the new entry. Neither the literal 42, the old maximum function, nor the
constant \code{contender\_5} occurs in \code{Hierarchy.lean} or
\code{Contender.lean}.

\begin{theorem}[Tower-height separation]\label{thm:towerentry}
Let $E_0(x)=x$ and $E_{j+1}(x)=2^{E_j(x)}$. Then
\begin{equation}\label{eq:towerfinal}
 E_{\Cfive+1}(64)<\Entry.
\end{equation}
In particular, the separation is far stronger than a fixed number of
successor steps or a single exponential in $\Cfive$.
\end{theorem}
\begin{proof}
Write $a_0=H_0(\leaf;\Boost)$ for the zeroth child of $\Theta$, and set
$m=R_{a_0}(64)$. By Theorem~\ref{thm:baseline}, $m>\Cfive$, hence
$m\ge\Cfive+1$. Corollary~\ref{cor:tower} applied to the root of $\Theta$
gives
\[
 E_m(64)<R_\Theta(64)=\Entry.
\]
For $x\ge1$, the sequence $E_j(x)$ is increasing: $E_1(x)=2^x>x$, and
successive exponentiations preserve the strict inequality. Therefore
$E_{\Cfive+1}(64)\le E_m(64)<\Entry$.
\end{proof}

The tower theorem uses only a small part of the grand index; it is a robust
lower bound, not an estimate of the entry's actual size. Most of the
higher-type and feedback construction is deliberately not unfolded by the
comparison proof.

\section{Why the construction is constructive rather than a loophole}
\label{sec:safety}

\subsection{No Busy Beaver and no unrestricted Lean evaluator}

Busy-Beaver-style maximization selects outputs of programs whose termination
cannot in general be determined by a total evaluator. Here every raw term
in each finite list already has a total denotation, established before the
maximum is defined. The evaluator does not run arbitrary Lean source,
search for consistency proofs, assume reflection for Lean itself, or
recognize a set of all total computable functions. The fixed-language escape
result is exactly why a universal self-maximizing interpretation would be
inappropriate.

\subsection{Infinite branching does not mean nontermination}

The only infinite object is a function giving children on demand. The tree
fold receives a function of recursive results and may inspect finitely many
children along a particular computation. A program supplied to the fold can
make the number of such inspections extraordinarily large, but its type and
recursor-based construction are what justify totality. There is no infinite
while-loop hidden behind a supposedly finite search bound.

\subsection{Avoiding an invalid oracle-monotonicity argument}

An object program can react nonmonotonically to its oracle inputs. Therefore
one should not infer, merely from $q\le q'$ pointwise, that every program's
output under $q$ is at most its output under $q'$. The proof strategy here
never needs such a statement. Envelope monotonicity is in the \emph{syntax
cutoff}, and stage comparisons use a specific term that calls the specified
child oracle. The tower theorem likewise uses a known program, not a global
monotonicity claim about interpreting every raw term.

\subsection{No unwarranted ordinal comparison}

The constructor $\Sup$ retains its branch order and syntactic presentation.
Two trees with equal externally assigned ordinal ranks may expose different
children at low numerical costs. Pointwise stage inequalities are proved
from those accessible children. An ordinal-ranking theorem, if developed,
would be additional work and would not replace the explicit coding
estimates needed for this contest.

\subsection{No claim of absolute optimality}

One can always define a successor or another reflection layer after a fixed
entry. The contest's computational time bounds and its prohibition of lazy
references shape a practical research problem, not a canonical greatest
natural number. Our proposal fixes a substantial language extension and
repeated higher-type amplification, and supplies a proof of a large margin.
It does not pretend that no more effective idea could be found with further
research.

\section{Lean artifact architecture}
\label{sec:implementation}

The definition path is deliberately separate from the baseline used in the
proof:
\[
\begin{gathered}
 \code{ListKit}\longrightarrow\code{Legacy}\longrightarrow\code{Language}
 \longrightarrow\code{Hierarchy}\longrightarrow\code{Contender},\\
 \code{Legacy}\longrightarrow\code{Baseline},\qquad
 (\code{Contender},\code{Baseline})\longrightarrow\code{Proofs}.
\end{gathered}
\]
\code{Audit} and \code{Smoke} are downstream validation modules.
\code{Legacy} supplies the old raw evaluator, not its winning constant;
\code{Baseline} alone supplies that constant.

\begin{center}
\begin{tabular}{p{.23\linewidth}p{.69\linewidth}}
\toprule
Module & Responsibility\\
\midrule
\code{ListKit.lean} & Elementary list operations, lookup, maximum, argmax.\\
\code{Legacy.lean} & Computational Coq-semantics port, weighted depth, raw enumerators.\\
\code{Language.lean} & Tree type, total extended interpreter, quotation, cumulative enumeration.\\
\code{Hierarchy.lean} & $R$, finite type levels, pointwise suprema, collapse, feedback, $\Theta$.\\
\code{Contender.lean} & The single proposed entry $R_\Theta(64)$.\\
\code{Baseline.lean} & The exact selected-term definition of $\Cfive$.\\
\code{Proofs.lean} & Occurrence lemmas, basic child and feedback inequalities, strict baseline proof.\\
\code{Audit.lean} & Transitive axiom queries for the entry and central theorems.\\
\code{Smoke.lean} & Tiny definitional examples, not giant-number evaluation.\\
\bottomrule
\end{tabular}
\end{center}

The maximum proof uses an explicit inductive list-occurrence predicate.
This avoids dependence on extensional finite-set APIs and makes the needed
membership derivations short. The final proof is intended to elaborate to
an application of the uniform baseline theorem, rather than normalize
$\Theta$ or run a maximum. The definition itself uses structural recursors
and does not employ a noncomputable semantic supremum.

The code's size is not evidence that the time limits are met. In particular,
Lean elaboration, generated code for recursors, imported theorem dependencies,
and conversion checking all require an actual toolchain run. These are
precisely what the acceptance script is designed to test.

\section{Validation performed and validation still required}
\label{sec:validation}

\subsection{Executed small-instance checks}

The Python model is an independent executable rendition of the raw
semantics, with explicit resource guards. It refuses large enumerations
rather than substituting bounded or approximate values into the mathematical
definition. It is not a Lean checker and does not certify totality.

The recorded run passed \textbf{62,065 assertions}. The old-language
enumeration results were:
\begin{center}
\begin{tabular}{rrrr}
\toprule
Depth & Raw terms & Types & Maximum numerical output\\
\midrule
1 & 2 & 1 & 0\\
2 & 10 & 2 & 1\\
3 & 146 & 5 & 2\\
4 & 24,976 & 26 & 3\\
\bottomrule
\end{tabular}
\end{center}
For the new language, take $q_0(i,n)=0$,
$q_1(i,n)=3i+n+7$, and $q_2(i,n)=\max(17-i-n,0)$:
\begin{center}
\begin{tabular}{rrrrrr}
\toprule
Depth & Raw terms & Types & $M_{q_0}$ & $M_{q_1}$ & $M_{q_2}$\\
\midrule
1 & 5 & 2 & 0 & 0 & 0\\
2 & 62 & 6 & 1 & 1 & 1\\
3 & 6,167 & 38 & 2 & 8 & 17\\
\bottomrule
\end{tabular}
\end{center}
The tiny results are not approximations to $\Cfive$ or $\Entry$.

The tests additionally cover argmax/maximum agreement, selection membership,
quotation and successor semantics, all enumerated depth bounds, cumulative
inclusion, absolute binder levels, malformed applications and casts, binary
numerals, the explicit depth-8 doubling and depth-14 exponentiation programs,
the depth-20 feedback witness at input 64, oracle-witness semantics, and
finite observations of higher-type branch identities. Binary inequalities
are checked through 9,999 in addition to their general mathematical proof.
The model tests power-of-two semantics only at small inputs and feedback
semantics only for small simple oracles; it never evaluates a giant stage.

\subsection{Unexecuted Lean validation}

There was no \code{lean}, \code{lake}, or \code{elan} executable in the
preparation container. Attempts to obtain a toolchain did not succeed.
The validation report therefore records
\begin{lstlisting}
"kernel_checked": false,
"timing_qualified": false,
"status": "not_run_compiler_unavailable"
\end{lstlisting}
A preliminary lexical source screen found no prohibited tokens in code after
stripping comments and strings. This is only a source screen. It cannot
establish elaboration, kernel acceptance, or the transitive axiom set of a
compiled theorem. The intended strict comparison proof contains no
placeholder, but unobserved syntax, API, termination-elaboration, simplifier,
or code-generation issues may remain.

\begin{center}
\begin{tabular}{p{.56\linewidth}p{.34\linewidth}}
\toprule
Claim or check & Delivered status\\
\midrule
Constructive mathematical definition & Fully specified\\
Mathematical proof $\Cfive<\Entry$ & Provided\\
Mathematical tower-height separation & Provided\\
Exact computational port argument & Provided; no cross-kernel certificate\\
Lean source for basic comparison & Provided; uncompiled\\
Lean source for tower/escape theorems & Not formalized in this package\\
Finite Python model & 62,065 assertions passed\\
Source placeholder screen & Passed, limited scope\\
Lean kernel and empty-axiom audit & Not run\\
15-second definition / 60-second proof & Unverified\\
Accepted contest submission & Not claimed\\
\bottomrule
\end{tabular}
\end{center}

\subsection{Reproduction and acceptance gates}

From the extracted archive root, the model and document are reproducible by
\begin{lstlisting}
python python/test_model.py
bash scripts/build_article.sh
\end{lstlisting}
On a machine with the pinned Lean toolchain installed, run
\begin{lstlisting}
python scripts/verify_lean.py
\end{lstlisting}
The script removes this package's generated Lean module artifacts, checks
its source screen, and compiles the five definition modules against the
prebuilt standard distribution under an aggregate 15-second budget. It then
compiles the baseline, proof, audit, and smoke modules under an aggregate
60-second budget. It records compiler output, elapsed times, and axiom
reports. Any compiler error, missing empty-axiom result, or time-budget
failure prevents an accepted status.

The required audit includes both the entry and its strict comparison:
\begin{lstlisting}
BigNumber.contender_6
BigNumber.contender_5_lt_contender_6
\end{lstlisting}
An empty dependency list is
required, not merely absence of a user-written \code{axiom} declaration.
The script's parser and timing policy are themselves ordinary validation
utilities, not a security-certified sandbox. A contest judge should also
review the source, compiler provenance, and machine description.

Do not run \code{\#eval contender\_6}, and do not attempt to print its decimal
expansion. The problem asks for a checked total definition and comparison,
not feasible normalization of the output.

\section{Further research questions}
\label{sec:future}

\paragraph{Complete and independently replay the strict formalization.}
The immediate concrete task is to run the pinned compiler, repair any
elaboration issues without weakening the mathematical specification, and
obtain empty axiom reports and disclosed timings. Rechecking with an
independent kernel implementation would strengthen the trust story. The
cross-language port should ultimately receive a machine-checked semantic
simulation rather than remain a paper argument plus finite tests.

\paragraph{Formalize the tower engine and the binary coding theorem.}
The strongest quantitative result is presently proved on paper. A systematic
library of context-aware quotation and depth certificates would make the
power-of-two and feedback witnesses straightforward to verify without
normalizing their outputs. Absolute binder levels make capture-avoiding
lifting an especially important component.

\paragraph{Optimize visible-child coding.}
The bound $2^{28}$ is deliberately conservative. Exact depth depends on both
bit length and population count in our binary encoding. Better addition
chains, shared descriptions, or a different explicitly charged numeral
language could improve it. The relevant optimization target is guaranteed
witness depth, not an informal textual abbreviation.

\paragraph{Replace extensional tree equalities by logical relations.}
Formalize the higher-type supremum/collapse identities in a way compatible
with the strict empty-axiom rule. A family of pointwise relations indexed by
finite type codes may prove observational equality without appealing to
functional extensionality as an axiom. One must also prove the evaluator
preserves that relation.

\paragraph{Determine genuine growth classifications.}
How does the finite-language envelope with W-recursion compare with named
fast-growing hierarchies under a specified coding? What are sharp lower and
upper bounds for the finite-type diagonal and feedback operation? An answer
requires precise ordinal notation or another semantic ranking together with
coding estimates. No identification with a famous ordinal is assumed here.

\paragraph{Investigate stronger constructive object languages.}
The upstream repository also includes work toward a System F evaluator
\cite{systemf}. Possible directions include a carefully stratified
polymorphic language, inductive families, or richer tree constructors, but
each extension must retain a total executable interpreter and satisfy the
actual definition-time gate. A larger semantic language is useful only when
its interpreter is genuinely specified and justified.

\paragraph{Separate semantic improvement from packaging tricks.}
A useful research standard would compare constructors, available recursion
principles, and budgeted embedding theorems, rather than reward hidden
precompiled definitions or unconstrained source abbreviation. The present
entry suggests measuring language inclusion and explicit stage simulation
alongside the traditional numerical comparison.

\section{Conclusion}

The proposed number is a closed constructive expression with a substantial
structural design: finite-language maximization, tree-indexed reflection,
arbitrary finite higher-type iteration, and feedback of semantic values
into type-level indices. Its relation to the actual upstream entry is not
left to growth-rate intuition. The successor-quotation witness proves
$\Cfive<\Entry$, and a separate depth-20 iteration witness proves the much
stronger inequality $E_{\Cfive+1}(64)<\Entry$.

The package makes the remaining boundary explicit. It provides mathematical
proofs, proposed Lean source, tests, and acceptance tooling. It does not
provide a successful Lean kernel run or a timing certificate. Thus it is a
serious research contender ready for the stated validation work, rather than
an already certified winner.

\appendix
\section{The central Lean definition, in ASCII notation}

The following is a readable ASCII rendering of the central definitions;
the actual source files are authoritative. It uses the interpreter and
finite envelope from \code{Language.lean}.
\begin{lstlisting}[basicstyle=\ttfamily\footnotesize]
def reflect : Tree -> Nat -> Nat :=
  Tree.rec
    (fun n => Language.envelope (fun _ _ => 0) (n + 4))
    (fun _ ih n =>
      Language.envelope (fun i x => ih i x) (n + 4))

def Level : Nat -> Type
  | 0 => Tree
  | n + 1 => Level n -> Level n

def supLevel : (k : Nat) -> (Nat -> Level k) -> Level k
  | 0, f => Tree.sup f
  | k + 1, f => fun x => supLevel k (fun n => f n x)

def omegaLevel (k : Nat) : Level (k + 2) :=
  fun f x => supLevel k (fun n => iterate f n x)

def seed (a : Tree) (s : Tree -> Tree) : (k : Nat) -> Level k
  | 0 => a
  | 1 => s
  | k + 2 => omegaLevel k

def collapse (a : Tree) (s : Tree -> Tree) :
    (k : Nat) -> Level k -> Tree
  | 0, x => x
  | k + 1, f => collapse a s k (f (seed a s k))

def high (a : Tree) (s : Tree -> Tree) (k : Nat) : Tree :=
  collapse a s (k + 2) (omegaLevel k)

def step (a : Tree) : Tree := Tree.sup (fun _ => a)

def boost (a : Tree) : Tree :=
  Tree.sup (fun
    | 0 => a
    | k + 1 => high a step (reflect a (k + 64)))

def grandIndex : Tree :=
  Tree.sup (fun k => high Tree.zero boost k)

def contender_6 : Nat := reflect grandIndex 64
\end{lstlisting}
Here \code{iterate} is ordinary natural-number recursion on the iteration
count. There is no call to the previous winning number in this definition.

\section{Provenance and reproducibility notes}

The upstream snapshot is pinned by:
\begin{description}[leftmargin=0pt,labelindent=0pt]
\item[Repository commit] \path{d56dce013e6264292533efdc5e1de71863c6fb86}.
\item[Contender.v blob] \path{f92db72deb0f3a4d4e3b575b4949eda73b1417a0}.
\item[README.md blob] \path{b2a2642114e50e86f4432f59aebf439617f78d18}.
\end{description}
The ported computational material retains attribution under the upstream
MIT license; its complete notice is in
\path{provenance/UPSTREAM_LICENSE.txt}. The package includes source metadata,
not a complete upstream repository mirror. No toolchain executable, font
file, or giant decimal numeral is included.

The article's test counts correspond to
\path{validation/python_tests.json}. The unverified Lean status is recorded
in \path{validation/lean_validation.json}. Re-running the scripts overwrites
these reports with the new observed results; such a run should be archived
with its machine details if used for contest adjudication.

\begin{thebibliography}{9}
\small\raggedright
\bibitem{repo}
Cody Roux and contributors, \emph{name-the-biggest-number}, source repository,
accessed 12 October 2026.
\url{https://github.com/codyroux/name-the-biggest-number}.

\bibitem{rules}
Cody Roux, \emph{name-the-biggest-number: README and contest rules}, pinned
snapshot described in Appendix B.
\url{https://github.com/codyroux/name-the-biggest-number/blob/d56dce013e6264292533efdc5e1de71863c6fb86/README.md}.

\bibitem{contender}
Cody Roux, Samuel Gr\"utter, and contributors, \emph{Contender.v}, especially
\code{cast\_impl}, \code{interp\_term}, \code{termsUpTo}, \code{maxBy}, and
\code{contender\_5}, pinned snapshot.
\url{https://github.com/codyroux/name-the-biggest-number/blob/d56dce013e6264292533efdc5e1de71863c6fb86/Contender.v}.

\bibitem{systemf}
Cody Roux, \emph{System\_F.v}, source in the same repository. Only the opening
normalization/evaluator development was inspected for this report.
\url{https://github.com/codyroux/name-the-biggest-number/blob/d56dce013e6264292533efdc5e1de71863c6fb86/System_F.v}.

\bibitem{leaninductive}
Lean contributors, \emph{The Lean Language Reference: Inductive Types},
accessed 12 October 2026.
\url{https://lean-lang.org/doc/reference/latest/The-Type-System/Inductive-Types/}.

\bibitem{leanrecursive}
Lean contributors, \emph{The Lean Language Reference: Recursive Definitions},
accessed 12 October 2026.
\url{https://lean-lang.org/doc/reference/latest/Definitions/Recursive-Definitions/}.

\bibitem{leanaxioms}
Jeremy Avigad, Leonardo de Moura, Soonho Kong, and Sebastian Ullrich,
\emph{Theorem Proving in Lean 4: Axioms and Computation},
accessed 12 October 2026.
\url{https://lean-lang.org/theorem_proving_in_lean4/Axioms-and-Computation/}.

\bibitem{leanrelease}
Lean contributors, \emph{Lean v4.35.0 release}, selected toolchain target;
not executed in the preparation environment.
\url{https://github.com/leanprover/lean4/releases/tag/v4.35.0}.

\bibitem{aaronson}
Scott Aaronson, \emph{Who Can Name the Bigger Number?}, background essay
linked by the contest repository.
\url{https://www.scottaaronson.com/writings/bignumbers.html}.
\end{thebibliography}
\end{document}
